\chapter{Das Siegelbandlogo}

Die verschiedenen Farbvarianten des Siegelbandlogos in \tubslatex werden
vom Paket \lstinline{tubslogo} zur Verfügung gestellt.
Die folgende Beschreibung bezieht sich auf den Funktionsumfang des Paketes.
Das Paket ist in allen verfügbaren Klassen bereits korrekt eingebunden.

Zur Verfügung gestellt werden eine Variante im rgb-Farbmodell, eine im
cmyk-Farbmodell, eine mit der Sonderfarbe HKS 15 oder (für Ausnahmefälle)
eine monochromen Variante.
Die Auwahl erfolgt dabei über Paketoptionen
(siehe Abschnitt \ref{options:color}).

\begin{figure}[!ht]
\begin{minipage}{0.5\textwidth}
  \centering
  \includegraphics[width=\tubslogoBaseWidth]{TUBraunschweig_RGB}
  \subcaption{Logo in Standardgröße (RGB)}
\end{minipage}
\begin{minipage}{0.5\textwidth}
  \centering
  \includegraphics[width=\tubslogoBaseWidth]{TUBraunschweig_4C}
  \subcaption{Logo in Standardgröße (CMYK)}
\end{minipage}

\begin{minipage}{0.5\textwidth}
  \centering
  \includegraphics[width=\tubslogoBaseWidth]{TUBraunschweig_SC}
  \subcaption{Logo in Standardgröße (HKS 15)}
\end{minipage}
\begin{minipage}{0.5\textwidth}
  \centering
  \includegraphics[width=\tubslogoBaseWidth]{TUBraunschweig_B}
  \subcaption{Logo in Standardgröße (schwarz)}
\end{minipage}
\caption{Verfügbare Farbversionen der Logos}
\end{figure}

\begin{Declaration}
  \Macro{tubslogo}\OParameter{Skalierung}\\
  \Macro{tubslogoAbs}\Parameter{Breite}
\end{Declaration}


Es werden 2 Befehle zum Einfügen des Logos in Dokumente bereit gestellt.
Mit dem Befehl \Macro{tubslogo} wird ein Logo in
standardkonformer Größe der verwendeten Papiergröße entsprechend skaliert
eingefügt. Die verwendete Papiergröße kann dabei als Paketoption übergeben
werden und ist standardmäßig auf \lstinline{a4paper} voreingestellt
(siehe Abschnitt \ref{options:papersize}).
Als optionales Argument kann auch ein frei gewählter Skalierungsfaktor
vorgegeben werden (siehe Abschnitt \ref{cmd:tubslogo}).

Mit \Macro{tubslogoAbs} kann über den Parameter \PName{Breite} dagegen
zusätzlich eine individuell gewählte absolute Breite angegeben werden.


\section{Paket-Optionen}

\subsection{Farbmodell}\label{options:color}

\begin{Declaration}
  \Option{rgb}\\
  \Option{cmyk}\\
  \Option{hsk}\\
  \Option{mono}
\end{Declaration}

Die Option \Option{rgb} bewirkt eine Darstellung des Logos in rgb-Farben.
Dies ist die Standardeinstellung.
Die Option \Option{cmyk} bewirkt eine Darstellung in cmyk-Druckfarben,
\Option{hks} wählt die Spezialfarbe \emph{HKS 15}.
Mit der Option \Option{mono} wird eine schwarz-weiß-Version des Logos geladen.

\subsection{Papiergröße}\label{options:papersize}

\begin{Declaration}
  \Option{a6paper}\\
  \Option{a5paper}\\
  \Option{a4paper}\\
  \Option{a3paper}\\
  \Option{a2paper}\\
  \Option{a1paper}\\
  \Option{a0paper}
\end{Declaration}
Bewirkt eine Skalierung des Logos auf dem Format entsprechende Größe.
Standardmäßig sind dies:

\begin{tabular}{lr}
  A6  & 60\% der Basisgröße   \\
  A5  & 70\% der Basisgröße   \\
  A4  & 100\% der Basisgröße  \\
  A3  & 140\% der Basisgröße  \\
  A2  & 200\% der Basisgröße  \\
  A1  & 280\% der Basisgröße  \\
  A0  & 400\% der Basisgröße   \\
\end{tabular}


\clearpage
\section{Befehle}

\begin{Declaration}
  \Macro{tubslogo}\OParameter{Breite rel.}
\end{Declaration}\label{cmd:tubslogo}

%   \item[\color{tuRed}\mdseries\ttfamily \textbackslash tubslogo%
%     {[\textcolor{tuGreenDark}{\sffamily\itshape $<$Breite rel.$>$}]}]
Erzeugt ein Logo, dass entsprechend des durch Angabe der Papiergröße
gewählten Skalierungsfaktors skaliert wird.

Wird das optionale Argument {\sffamily\itshape $<$Breite rel.$>$}
verwendet, so kann ein freier Skalierungsfaktor ausgehend von der
Standardgröße für DINA4-Dokumente ($70$mm$\times 26$mm) übergeben werden.

\lstinline!\tubslogo[0.6]! erzeugt beispielsweise ein auf 60\% der 
Standardgröße für DINA4-Dokumente skaliertes Logo.

\begin{Declaration}
  \Macro{tubslogoAbs}\PParameter{Breite abs.}
\end{Declaration}\label{cmd:tubslogo}
Dieser Befehl erlaubt die Erzeugung eines Logos mit absolut vorgegebener
Breite.\footnote{Generell ist die Verwendung von standardkonformen Größen
einer individuellen Skalierung vorzuziehen!}

\begin{example}
  \lstinline!\tubslogoAbs{5cm}! erzeugt ein 5cm breites Logo.
\end{example}


\section{Längen}

\begin{Declaration}
  \Length{tubslogoBaseWidth}\\
  \Length{tubslogoBaseHeight}
\end{Declaration}
Basisbreite (70mm) und -höhe des Logos (26mm).
Dies entsprechen der korrekten Darstellung für das \mbox{DIN\,A4}-Format.
Die Größen für die weiteren DIN-Formate leiten sich daraus über
Skalierungsfaktoren ab.

\begin{Declaration}
  \Length{tubslogoWidth}\\
  \Length{tubslogoHeight}
\end{Declaration}
Breite bzw. Höhe des automatisch skalierten Logos.
Wird automatisch anhand der angegebenen Papiergröße berechnet und
kann zur dynamischen Positionierung verwendet werden.
